\documentclass[12pt]{article}
\usepackage[margin=1in]{geometry}
\usepackage{amsmath}
\usepackage{amssymb}
\usepackage{amsthm}
\usepackage{enumitem}
\usepackage{xcolor}
\usepackage{pagecolor}
\usepackage{marginnote}
\usepackage{fancyhdr}
\usepackage{bm}

% Define cream background color
\definecolor{cream}{HTML}{faf7f1}
\pagecolor{cream}

% Define defn and defeq commands
\newcommand{\defn}[1]{\textbf{#1}}
\newcommand{\defeq}{\mathrel{\vcenter{\baselineskip0.5ex \lineskiplimit0pt
                     \hbox{\scriptsize.}\hbox{\scriptsize.}}}=}
\newcommand{\T}{\top}
\newcommand{\F}{\bot}

\newtheoremstyle{proofstyle}
  {}{}{\itshape}{}{\bfseries}{.}{ }{}
\theoremstyle{proofstyle}
\newtheorem*{proof*}{Proof}

% Define sidenote command (simplified - uses marginnote)
\newcommand{\sidenote}[1]{\marginnote{\small\itshape #1}}

% Title formatting
\makeatletter
\renewcommand{\maketitle}{%
  \begin{flushleft}
    {\Large\bfseries CS173: Discrete Structures}\\[0.3em]
    {\Large\bfseries Problem Set 1}\\[0.5em]
    {\normalsize Mohammed Al Salhi}
  \end{flushleft}
  \vspace{1em}
}
\makeatother

\AtBeginDocument{\mathversion{bold}}

\begin{document}

\maketitle

\begin{enumerate}[left=0pt, itemsep=2em, parsep=1em]

\item Prove each of the following statements without truth tables.

\begin{enumerate}[label=(\alph*), itemsep=1.5em, leftmargin=2em]

\item Prove $p \to p$ is a tautology for any proposition $p$.

\begin{proof*}
Let $p$ be a proposition.
\begin{align*}
p \to p &\equiv \neg p \vee p && \text{(By Conditional Disintegration)} \\
&\equiv \T && \text{(By Complement)}
\end{align*}
$\therefore$ $p \to p$ is a tautology. \qed
\end{proof*}

\item Prove $p \to q \equiv \neg q \to \neg p$ for all propositions $p, q$.

\begin{proof*}
Let $p$ and also $q$ be propositions.
\begin{align*}
p \to q &\equiv \neg p \vee q && \text{(By Conditional Disintegration)} \\
&\equiv q \vee \neg p && \text{(By Commutativity)} \\
&\equiv \neg (\neg q ) \vee \neg p && \text{(By Double Negation)} \\
&\equiv \neg q \to \neg p && \text{(By Conditional Disintegration)}
\end{align*}
$\therefore p \to q \equiv \neg q \to \neg p$ for all propositions $p, q$. \qed
\end{proof*}

\item Prove $\neg(p \to q) \equiv p \wedge \neg q$ for all propositions $p, q$.

\begin{proof*}
Let $p$ and also $q$ be propositions.
\begin{align*}
\neg(p \to q) &\equiv \neg(\neg p \vee q) && \text{(By Conditional Disintegration)} \\
&\equiv \neg(\neg p) \wedge \neg q && \text{(By De Morgan's Law)} \\
&\equiv p \wedge \neg q && \text{(By Double negation)}
\end{align*}
$\therefore \neg(p \to q) \equiv p \wedge \neg q$ for all propositions $p$ and also $q$. \qed
\end{proof*}

\item
\begin{samepage}
Prove $(p \wedge (p \to q)) \to q$ is a tautology for all propositions $p, q$.

\begin{proof*}
Let $p$ and also $q$ be propositions.
\begin{align*}
(p \wedge (p \to q)) \to q &\equiv (p \wedge (\neg p \vee q)) \to q && \text{(By Conditional Disintegration)} \\
&\equiv ((p \wedge \neg p) \vee (p \wedge q)) \to q && \text{(By Distributivity)} \\
&\equiv ((\neg p \wedge p) \vee (p \wedge q)) \to q && \text{(By Commutativity)} \\
&\equiv (\F \vee (p \wedge q)) \to q && \text{(By Complement)} \\
&\equiv (p \wedge q) \to q && \text{(By Identity)} \\
&\equiv \neg(p \wedge q) \vee q && \text{(By Conditional Disintegration)} \\
&\equiv (\neg p \vee \neg q) \vee q && \text{(By De Morgan's law)} \\
&\equiv \neg p \vee (\neg q \vee q) && \text{(By Associativity)} \\
&\equiv \neg p \vee \T && \text{(By Complement)} \\
&\equiv \T && \text{(By Domination)}
\end{align*}
$\therefore (p \wedge (p \to q)) \to q$ is a tautology for all propositions $p, q$. \qed
\end{proof*}
\end{samepage}

\item Prove $(p \vee q) \to r \equiv (p \to r) \wedge (q \to r)$ for all propositions $p, q, r$.

\begin{proof*}
Let $p, q,$ and also $r$ be propositions.
\begin{align*}
(p \vee q) \to r &\equiv \neg (p \vee q) \vee r && \text{(By Conditional Disintegration)} \\
&\equiv (\neg p \wedge \neg q) \vee r && \text{(By De Morgan's law)}\\
&\equiv r \vee (\neg p \wedge \neg q) && \text{(By Commutativity)}\\
&\equiv (r \vee \neg p) \wedge ( r \vee \neg q) && \text{(By Distributivity)}\\
&\equiv (\neg p \vee r) \wedge (\neg q \vee r) && \text{(By Commutativity)}\\
&\equiv (p \to r) \wedge (q \to r) && \text{(By Conditional Disintegration)}
\end{align*}
$\therefore (p \vee q) \to r \equiv (p \to r) \wedge (q \to r)$ for all propositions $p, q, r$. \qed
\end{proof*}

\end{enumerate}

\item Prove each of the following statements for all propositions $\varphi, \psi, \xi$.

\begin{enumerate}[label=(\alph*), itemsep=1.5em, leftmargin=2em]

\item $(\varphi \to \psi), (\psi \to \xi) \vdash \varphi \to \xi$ \hfill Hypothetical Syllogism

\begin{proof*}
    Let $\varphi, \psi, \xi$ be arbitrary propositions, and suppose $(\varphi \to \psi)$ and $(\psi \to \xi)$.
    
    To prove $(\varphi \to \xi)$, by the deduction rule, it suffices to show $\varphi \vdash \xi$.
    
    \begin{quote}
        Assume $\varphi$. \\
        By the assumption $(\varphi \to \psi)$, we have $\psi$ by modus ponens. \\
        By the assumption $(\psi \to \xi)$, we have $\xi$ by modus ponens. \\
        Thus $\varphi \vdash \xi$.
    \end{quote}
    
    Therefore, by the deduction rule, $(\varphi \to \xi)$. \qed
\end{proof*}

\item $\vdash \varphi \to \varphi$

\begin{proof*}
    Let $\varphi$ be an arbitrary proposition.
    
    To prove $(\varphi \to \varphi)$, by the deduction rule, it suffices to show $\varphi \vdash \varphi$.
    
    \begin{quote}
        Assume $\varphi$. \\
        Then we have $\varphi$.
    \end{quote}
    
    Therefore, by the deduction rule, $(\varphi \to \varphi)$. \qed
\end{proof*}

\item $\vdash \top$

\begin{proof*}
    From part (b), we have $\varphi \to \varphi$ for any proposition $\varphi$.
    
    From part 1(a) $\varphi \to \varphi \equiv \T $
    
    Therefore, we have $\top$. \qed
\end{proof*}


\item $\varphi, \psi \vdash \varphi \wedge \psi$ \hfill Conjunction Introduction

\begin{proof*}
    Let $\varphi$ and $\psi$ be a arbitrary propositions. Assume $\varphi$ and also assume $\psi$.
    
    To prove $\varphi \land \psi$, we use reductio ad absurdum. 
    
    We will show $(\neg(\varphi \land \psi) \vdash \psi)$ and $(\neg(\varphi \land \psi) \vdash \neg\psi)$.
    
    Assume $\neg(\varphi \land \psi)$.
    \begin{align*}
        \neg(\varphi \land \psi) &\equiv \neg\varphi \lor \neg\psi && \text{(By De Morgan's law)} \\
        &\equiv \varphi \to \neg\psi && \text{(By Conditional Disintegration)}
    \end{align*}
        We now have $(\varphi \to \neg\psi)$ and $\varphi$ from the intial assumption \\
        Thus we have $\neg\psi$ by Modus Ponens. However, we originally assumed $\psi$, so there is a contradiction. \\    
    Therefore, by reductio ad absurdum, we conclude $\varphi \land \psi$.\qed
\end{proof*}

\newpage
\item $\varphi \wedge \psi \vdash \varphi$ \hfill Conjunction Elimination

\begin{proof*}
    Let $\varphi$ and $\psi$ be arbitrary propositions. Assume $\varphi \land \psi$.
    
    To prove $\varphi$ we use reductio ad absurdum. Assume $\neg\varphi$.
    
    By conjunction introduction we have $(\varphi \land \psi) \land \neg\varphi$:
    \begin{align*}
        (\varphi \land \psi) \land \neg\varphi &\equiv \varphi \land (\psi \land \neg\varphi) && \text{(By Associativity)} \\
        &\equiv \varphi \land (\neg\varphi \land \psi) && \text{(By Commutativity)} \\
        &\equiv (\varphi \land \neg\varphi) \land \psi && \text{(By Associativity)} \\
        &\equiv (\neg \varphi \land \varphi) \land \psi && \text{(By Commutativity)} \\
        &\equiv \bot \land \psi && \text{(By Complement)} \\
        &\equiv \bot && \text{(By Domination)}
    \end{align*}
    
    Thus we have $\bot$.
    
    However, from part 2(c), we have $\top$.
    
    Furthermore:
    \begin{align*}
        \top &\equiv \neg \varphi \lor \varphi && \text{(By Complement)} \\
        &\equiv \neg(\neg(\neg\varphi) \land \neg\varphi) && \text{(By De Morgan's law)} \\
        &\equiv \neg(\varphi \land \neg\varphi) && \text{(By Double negation)} \\
        &\equiv \neg(\neg\varphi \land \varphi) && \text{(By Commutativity)} \\
        &\equiv \neg\bot && \text{(By Complement)}
    \end{align*}
    We have both $\bot$ and $\neg\bot$, which is a contradiction.
    
    Therefore, by reductio ad absurdum, $\varphi$ holds. \qed
\end{proof*}

\end{enumerate}

\end{enumerate}

\end{document}